\documentclass[11pt]{article}

\usepackage{geometry}
\usepackage{setspace}
\usepackage{enumitem}
\usepackage{array}
\usepackage[colorlinks=true, urlcolor=black]{hyperref}

\newcolumntype{L}[1]{>{\raggedright\let\newline\\\arraybackslash\hspace{0pt}}m{#1}}

\geometry{left=0.7in,right=0.7in,top=0.8in,bottom=0.8in}
\setlength{\parindent}{0in}

% \setstretch{1.2}

\pagestyle{empty}

\begin{document}


\begin{center}

\section*{Keqing Liu}

\end{center}

\begin{center}
Associate Professor, Wang Yanan Institute for Studies in Economics,  Xiamen University \\ 
\textbf{Email}: liu.keqing@outlook.com \quad \textbf{Homepage}: \url{https://keqing-liu.github.io/} \\
\vspace{0.1in}
\textbf{Phone No.}: (+86) 13771110112 \quad  \textbf{Citizenship}: Chinese | Canadian PR: In Process

\end{center}

\section*{Professional Summary}
Quantitative Macro Economist and Data Scientist with extensive expertise in building predictive and structural macroeconomic models (DSGE) alongside machine learning algorithms (e.g., Random Forest).
Proven ability to leverage high-dimensional datasets to simulate how agents (households, firms, and governments) look ahead and react to strategic shocks.
Adept at running scalable counterfactual simulations to forecast economic trends and evaluate policy or business impacts—such as quantifying the systemic effects of interest rate hikes or geopolitical tensions.
Skilled in Python and MATLAB, with a strong record of translating complex quantitative analysis into clear, strategic recommendations for both technical and non-technical audiences.
\section*{Technical Skills}
\begin{itemize}
    \item \textbf{Quantitative Modeling}: DSGE models, statistical modeling, econometric estimation, time series analysis (e.g., ARIMA, SVAR), machine learning (e.g., Random Forest),  forecasting
    \item \textbf{Programming}: Python (NumPy, Pandas, SciPy, Scikit-learn), SQL, MATLAB, R, Stata, Git
    \item \textbf{Data Analysis}: Regression analysis, data visualization, counterfactual analysis, causal inference
\end{itemize}

\section*{Selected Technical Projects}
\begin{itemize}
    \item \textbf{Personal Quantitative Research System:} Built a Python and SQLite-based workflow to organize market data, backtest trading strategies, evaluate entry and exit signals, and support disciplined investment analysis.
\end{itemize}


\section*{Education}

\begin{tabular}{L{2.4cm} L{5cm} L{8cm}}
2013 - 2017 &  Ph.D. in Economics & University of Exeter (ESRC scholarship)\\
2012 - 2013 & MSc Economics (Distinction) & University of Exeter\\
2008 - 2012 & B.A. in Economics &  Huazhong University of Science \& Technology 
\end{tabular}

\section*{Work Experience}

\begin{tabular}{L{3.2cm} L{3.5cm} L{9cm}}
2025.09 - Present & Associate Professor & Xiamen University \\
2018.09 - 2025.08 & Assistant Professor & Xiamen University \\
2017.08 - 2018.07 & Associate Professor & Shanghai University of Finance and Economics \\
\end{tabular}

\section*{Research Expertise}

I specialize in integrating macro-financial economic theory with advanced data science techniques to solve complex forecasting, risk assessment, and policy evaluation problems. My structural and empirical expertise includes:

\begin{itemize}
    \item \textbf{Economic Modeling \& Policy Simulation:} Expert in developing and estimating Dynamic Stochastic General Equilibrium (DSGE) models with frictional banking sectors. Adept at applying nonlinear and perturbation solution methods to run high-dimensional ``what-if'' policy experiments, simulating and forecasting the business cycle impacts of regulatory or interest rate shocks under diverse economic scenarios.
    \item \textbf{Machine Learning \& Risk Classification:} Applied Random Forest methods to bank-level balance sheet and regional macroeconomic data to study financial risk patterns. This work combined micro-level bank characteristics with broader economic conditions to support risk classification, early-warning analysis, and credit risk assessment.
    \item \textbf{Time-Series Estimation \& Econometrics:} Proficient in advanced statistical modeling and econometric methods (e.g., ARIMA, SVAR). Developed integrated New Keynesian frameworks that combine structural micro-foundations with time-series estimation techniques to solve optimal policy and time-inconsistency problems.
\end{itemize}

\section*{Policy Experience}

I have direct experience translating research into policy insights.
During my tenure at Shanghai University of Finance and Economics, I also served as a Senior Research Associate, contributing to a team that produced quarterly policy research reports on China's economy for the central bank and government institutions, providing analysis intended to inform policy discussions and strategic assessment.
My focus within the group was on monetary policy analysis.

\section*{Selected Publications}

\begin{enumerate}[label={[\arabic*]}]
\item{Habit formation and news-driven business cycles (2025), \textit{Review of Economic Dynamics}, 56, 101273,  with Fengqi Liu and Jianpo Xue.}
\item{Should macroprudential policy be countercyclical? (2024), \textit{Journal of Economic Dynamics and Control}, 158, 104765,  with Yoske Igarashi.}
\item {Credit expansion, bank liberalization, and structural change in bank asset accounts (2021), \textit{Journal of Economic Dynamics and Control}, 124, 104066, with Michael Fan.} 
\end{enumerate}


\section*{Selected Conference Presentations}
\begin{enumerate}[label={[\arabic*]}]
\item { Discussant. Economic Cycles and Growth Conference, in Xiamen, 2024: `Uncovering Frictions to Long-Term Borrowing'.}
\item { Asian Meeting of the Econometric Society (AMES) in Xiamen, 2019: `Credit Expansion, Bank Liberalization, and Structural Change in Bank Asset Accounts'.}
\item { Money Macro Finance Annual Conference (MMF) in Bath, 2016: `Bank Equity and Macroprudential Policy'.}
\end{enumerate}


\section*{Teaching (in English)}
\begin{tabular}{L{11cm} L{4cm}}
Money and Banking & 2018 - 2025 \\
Advanced Macroeconomics: Money, Banking, and Finance & 2023 - 2025 \\
Advanced Macroeconomics II & 2019 - 2022 \\
Intermediate Microeconomics (SUFE) & 2017 - 2018 \\
\end{tabular}

\end{document}
